\section{研究局限与展望}

\subsection{研究局限}

样本覆盖与代表性存在不足。本研究仅选取杭甬温三地开展调研，未覆盖浙江其余地市。同时本次调查样本以中青年、高学历群体为主，老年、低龄等敏感人群样本占比偏低，可能使噪声感知结果存在一定偏差，难以完整反映全年龄段居民的真实声环境体验。

数据采集存在时空局限。问卷调查为截面数据，仅反映特定时段居民噪声感知，未开展长期追踪调查，无法揭示噪声感知与治理成效的动态变化；监测数据依赖官方公开平台，部分区域数据完整性不足。

调研实施与传播条件存在局限。本次调研由大学生团队开展，人力、时间、传播渠道等资源均较为有限，且线上问卷主要依托社交网络传播，进一步加剧了老年群体与未成年群体的样本缺失问题，同时也对全面收集基层真实诉求造成了影响。

研究方法存在一定局限性。有序Logistic回归仅验证核心影响因素，未纳入经济水平、城市建设密度等宏观变量；质性访谈样本量较少，对噪声治理深层矛盾的挖掘不够深入。

治理建议的实操性有待细化。本研究提出的优化建议偏向宏观层面，针对不同城市、不同噪声源的具体实施方案、成本测算、效果评估等内容未展开细化，落地执行的细节支撑不足。

\subsection{研究展望}

未来调研可覆盖浙江全部11个地市，线上线下同步发放问卷，为老年、低龄群体提供代填协助，均衡各分层样本结构，提升结论全省普适性。

联合社区、高校、公益组织开展协同调研，拓宽线下推广与入户调查渠道，弥补大学生团队资源不足的短板；采用面板数据进行年度追踪调查，结合噪声监测、投诉数据与居民感知数据，长期跟踪噪声污染演变、治理成效与居民满意度变化，挖掘动态影响机制。

引入地理信息系统（GIS）、空间计量模型等方法，结合城市规划、经济发展等宏观数据，开展噪声污染空间分布与影响因素精细化研究，针对不同城市和不同噪声源制定可量化、可操作、可评估的治理方案，结合成本与效益分析测算治理投入与民生收益，为浙江省“浙里宁静”建设、宁静城市创建提供更具实操性的决策支撑。

探索跨学科融合，创新治理模式，融合环境科学、社会学、心理学、公共管理等多学科理论，构建噪声治理理论框架，为全国城市化进程中的噪声污染防治提供浙江经验。

